\documentclass[11pt,a4paper]{report}

% --- Packages indispensables ---
\usepackage[utf8]{inputenc}
\usepackage[T1]{fontenc}
\usepackage[french]{babel}
\usepackage{geometry}
\geometry{top=2.5cm, bottom=2.5cm, left=2.5cm, right=2.5cm}
\usepackage{amsmath, amsfonts, amssymb}
\usepackage{graphicx}
\usepackage{booktabs}
\usepackage{array}
\usepackage{multirow}
\usepackage{xcolor}
\usepackage{listings}
\usepackage{fancyhdr}
\usepackage{hyperref}
\usepackage{tcolorbox}
\usepackage{float}
\usepackage{microtype}
\usepackage{titlesec}

% --- Couleurs personnalisées ---
\definecolor{primary}{RGB}{18, 69, 89}      % Bleu marine sombre
\definecolor{secondary}{RGB}{89, 131, 146}  % Bleu sarcelle
\definecolor{accent}{RGB}{40, 167, 69}      % Vert agricole
\definecolor{warn}{RGB}{217, 83, 79}        % Rouge alerte
\definecolor{codebg}{RGB}{245, 247, 250}    % Gris clair code
\definecolor{codecomment}{RGB}{108, 117, 125}

% --- Configuration de hyperref ---
\hypersetup{
    colorlinks=true,
    linkcolor=primary,
    citecolor=secondary,
    urlcolor=accent,
    pdftitle={Rapport de Projet Edge AI - Irrigation Intelligente Souss-Massa},
    pdfauthor={Équipe Projet Edge AI},
    pdfsubject={Systèmes Embarqués, TinyML et Fog Computing}
}

% --- Configuration de listings (Code C / Python / JSON) ---
\lstset{
    backgroundcolor=\color{codebg},
    basicstyle=\ttfamily\footnotesize,
    breaklines=true,
    captionpos=b,
    commentstyle=\color{codecomment}\itshape,
    keywordstyle=\color{primary}\bfseries,
    stringstyle=\color{accent},
    numbers=left,
    numberstyle=\tiny\color{codecomment},
    numbersep=7pt,
    frame=single,
    rulecolor=\color{secondary!40},
    tabsize=2,
    showstringspaces=false
}

% --- Entêtes et pieds de page ---
\pagestyle{fancy}
\fancyhf{}
\fancyhead[L]{\small\color{secondary}\textbf{Projet Edge AI} -- Irrigation Intelligente}
\fancyhead[R]{\small\color{secondary}Région du Souss-Massa}
\fancyfoot[L]{\small\color{codecomment}Dépôt : \url{https://github.com/Nouzha-hs/projet-edge-ai.git}}
\fancyfoot[R]{\thepage}
\renewcommand{\headrulewidth}{0.4pt}
\renewcommand{\footrulewidth}{0.4pt}

% =========================================================================
\begin{document}

% --- Page de Garde ---
\begin{titlepage}
    \centering
    \vspace*{1cm}
    
    {\Huge \textbf{\color{primary}SYSTÈME D'IRRIGATION INTELLIGENTE AUTONOME BASÉ SUR L'EDGE AI (TINYML)}}\\[0.5cm]
    {\LARGE \textbf{\color{secondary}Application Éco-Responsable à l'Agriculture Aride du Souss-Massa}}\\[1.5cm]
    
    \rule{\linewidth}{1.5pt}\\[0.4cm]
    {\Large \textbf{RAPPORT TECHNIQUE D'INGÉNIERIE COMPLET}}\\[0.2cm]
    \rule{\linewidth}{1.5pt}\\[2cm]
    
    \begin{minipage}{0.48\textwidth}
        \textbf{\large Auteurs \& Rôles :}\\[0.3cm]
        \textbf{Membre 1 :} Ingénierie Embarquée \& Edge AI (STM32)\\
        \textbf{Membre 2 :} Data Science \& Climatologie (Copernicus / IA)\\
        \textbf{Membre 3 :} Fog Computing, IoT \& Supervision (Node-RED)
    \end{minipage}
    \hfill
    \begin{minipage}{0.48\textwidth}
        \flushright
        \textbf{\large Cadre Technique :}\\[0.3cm]
        \textbf{Plateforme :} STM32 Nucleo-F401RE\\
        \textbf{Cœur :} ARM Cortex-M4 @ 84 MHz\\
        \textbf{Supervision :} MQTT, Node-RED, Python\\
        \textbf{Données :} Copernicus ERA5 (731 jours)\\
        \textbf{Date :} Novembre 2026
    \end{minipage}
    
    \vfill
    
    \begin{tcolorbox}[colback=primary!5,colframe=primary,arc=3mm]
        \centering
        \textbf{Dépôt de Code Source & Démonstrateur GitHub :} \\
        \href{https://github.com/Nouzha-hs/projet-edge-ai.git}{\ttfamily https://github.com/Nouzha-hs/projet-edge-ai.git}
    \end{tcolorbox}
    
    \vspace*{0.5cm}
\end{titlepage}

% --- Résumé et Abstract ---
\chapter*{Résumé Exécutif}
\addcontentsline{toc}{chapter}{Résumé Exécutif}

Le présent rapport documente la conception, l'optimisation et la validation d'un système cyber-physique complet d'irrigation de précision basé sur l'**Edge AI (TinyML)** déployé sur cible **STM32 Nucleo-F401RE** (ARM Cortex-M4 à 84 MHz) pour la plaine du **Souss-Massa (Taroudant/Agadir, Maroc : $30.42^\circ\text{N}, -8.87^\circ\text{W}$)**.

Face aux contraintes climatiques arides, aux épisodes caniculaires de \textit{Chergui} et à la raréfaction critique des nappes phréatiques, ce projet résout le double problème des minuteries conventionnelles (gaspillage d'eau aveugle) et des solutions Cloud distantes (zones blanches sans réseau, coût et latence).

Les trois modules du projet réalisés en synergie par l'équipe sont :
\begin{enumerate}
    \item \textbf{Module Edge / Embarqué (Membre 1) :} Synthèse d'un modèle d'arbre de décision directement en langage C ANSI sans allocation dynamique (`malloc`), affichant une empreinte de **100 octets en Flash**, **0 octet de RAM dynamique** et une latence d'inférence de **0.183 $\mu$s** (15.4 cycles CPU). Intégration d'un coupe-circuit matériel de sécurité (\textit{Watchdog} 30 secondes) garantissant un risque d'inondation nul (0.0\%).
    \item \textbf{Module Data Science \& TinyML (Membre 2) :} Traitement de **731 journées de relevés réels Copernicus ERA5** (2023--2024), formulation de l'évapotranspiration de référence ($ET_0$) selon le standard international \textbf{FAO-56 Penman-Monteith}, et benchmarking multi-modèles (Arbre de Décision : 96.72\% de précision, Forêt Aléatoire : 97.27\%, Réseau MLP : 97.27\%).
    \item \textbf{Module Fog Computing \& Supervision (Membre 3) :} Architecture distribuée avec passerelle série-MQTT (\textit{Eclipse Mosquitto}), simulateur environnemental réaliste du Souss, tableau de bord \textit{Node-RED 2.0} et benchmark écologique démontrant une économie d'eau de **15.0\% à 32.3\%** par rapport aux minuteries à seuil fixe.
\end{enumerate}

\vspace{0.5cm}

\begin{otherlanguage}{english}
\section*{Abstract}
This technical report details the full engineering, design, and deployment of an autonomous, low-power smart irrigation system leveraging **Edge AI (TinyML)** on an **STM32 Nucleo-F401RE** microcontroller (ARM Cortex-M4 @ 84 MHz) targeted at the arid **Souss-Massa** plain (Morocco: $30.42^\circ\text{N}, -8.87^\circ\text{W}$). 

Driven by 731 days of real **Copernicus ERA5** climate observations (2023--2024) and the **FAO-56 Penman-Monteith** reference evapotranspiration ($ET_0$) metric, an ultra-compact ANSI C Decision Tree classifier was compiled with zero dynamic memory allocation (`malloc`), requiring only **100 bytes of Flash**, **0 bytes of dynamic RAM**, and executing within **0.183 $\mu$s** (15.4 clock cycles). A hardware-grade 30-second failsafe watchdog eliminates flooding risks under sensor failure. 

At the Fog layer, a local **Eclipse Mosquitto MQTT broker**, a serial-MQTT bridge, and a **Node-RED 2.0 telemetry dashboard** deliver real-time supervisory control. Multi-scenario comparative benchmarks reveal annual water savings between **15.0\% and 32.3\%** (up to 1,391 L/year saved per zone) over conventional fixed-schedule timers, demonstrating high environmental and operational impact for precision agriculture.
\end{otherlanguage}

\tableofcontents

% =========================================================================
\chapter{Introduction \& Contexte Régional du Souss-Massa}

\section{La Problématique Hydrique dans le Souss-Massa}
La plaine du Souss-Massa constitue la locomotive de la production maraîchère et agrumicole marocaine. Cependant, elle se situe dans un domaine bioclimatique aride à semi-aride caractérisé par :
\begin{itemize}
    \item Une pluviométrie moyenne annuelle très faible (inférieure à 200 mm/an), concentrée sur de rares épisodes orageux hivernaux.
    \item Un déficit d'alimentation chronique des nappes phréatiques (Chtouka, nappe du Souss), provoquant un abaissement du niveau piézométrique de plusieurs mètres par an.
    \item Le phénomène météorologique du \textbf{Chergui} : vent continental d'est/sud-est en provenance du Sahara, caractérisé par des températures extrêmes ($> 40^\circ\text{C}$), une humidité relative chutant en deçà de 15\%, et une demande évaporative journalière ($ET_0$) dépassant 9 mm/jour.
\end{itemize}

\section{Limites des Systèmes d'Irrigation Conventionnels}
Actuellement, les exploitations agricoles de la région utilisent principalement des électrovannes asservies à des minuteries programmables à seuil fixe ou à calendrier fixe. Cette approche présente des défaillances majeures :
\begin{enumerate}
    \item \textbf{Irrigation Aveugle :} L'arrosage s'enclenche de façon rigide même si le sol a bénéficié de précipitations récentes ou si l'air est saturé en humidité, gaspillant ainsi une ressource inestimable.
    \item \textbf{Incapacité d'Anticipation :} Lors des canicules soudaines ou des journées de Chergui, le débit prédéfini est insuffisant, entraînant un stress hydrique aigu pour les cultures.
    \item \textbf{Risque Catastrophique de Noyade Racinaire :} En cas de rupture d'un tuyau ou de panne de capteur analogique bloqué à l'état sec, la pompe fonctionne en continu sans surveillance, provoquant l'asphyxie des racines.
\end{enumerate}

\section{Pourquoi l'Edge AI Frugale plutôt que le Cloud ?}
Bien que des solutions d'irrigation connectées au Cloud existent sur le marché, elles s'avèrent inadaptées à la réalité opérationnelle des fermes du Souss :
\begin{itemize}
    \item \textbf{Zones Blanches et Ruptures de Réseau :} Dans les zones rurales reculées, les réseaux 4G et Wi-Fi sont instables. Un système dépendant du Cloud cesse d'irriguer en cas de déconnexion.
    \item \textbf{Coûts Récurrents :} Les abonnements cellulaires et de serveurs Cloud constituent un fardeau financier pour les petits exploitants.
    \item \textbf{Efficacité Énergétique :} L'émission radio permanente vers des stations cellulaires draine rapidement les batteries.
\end{itemize}

La démarche adoptée dans ce projet consiste à déporter l'intelligence artificielle directement à la périphérie (\textbf{Edge AI}), à l'intérieur du microcontrôleur assurant le contrôle physique de la motopompe.

% =========================================================================
\chapter{Architecture Cyber-Physique Globale}

Le système est articulé autour d'une architecture modulaire à trois niveaux :
\begin{figure}[H]
    \centering
    \begin{tcolorbox}[colback=white,colframe=primary,arc=2mm,title=\textbf{Architecture Distribuée Edge-Fog-Supervision}]
    \begin{verbatim}
+-----------------------------------------------------------------------+
|                NIVEAU 3 : SUPERVISION & DASHBOARD (NODE-RED)          |
|  - Dashboard interactif (Jauges T°/Humidité, Courbes ET0, État Pompe) |
|  - Commande manuelle dérogatoire & Archivage des décisions (CSV)      |
+-----------------------------------^-----------------------------------+
                                    | WebSockets / HTTP (Port 1880)
+-----------------------------------v-----------------------------------+
|                NIVEAU 2 : FOG COMPUTING & PASSERELLE LOCALE           |
|  - Broker MQTT local (Eclipse Mosquitto, Port 1883)                   |
|  - Pont Série <-> MQTT (bridge_serial_mqtt.py)                        |
|  - Simulateur environnemental réaliste du Souss (simulateur_souss.py) |
+-----------------------------------^-----------------------------------+
                                    | Liaison Série UART (115200 bauds)
+-----------------------------------v-----------------------------------+
|                NIVEAU 1 : EDGE AI (STM32 NUCLEO-F401RE / CORTEX-M4)   |
|  - Acquisition capteur capacitif de sol (ADC1 12 bits sur PA0)        |
|  - Inférence locale en C pur (Arbre de Décision TinyML, 0.183 us)     |
|  - Pilote actionneur relais de pompe (GPIO Output sur PA5)            |
|  - Coupe-circuit de sécurité Watchdog (30 secondes maximum)           |
+-----------------------------------------------------------------------+
    \end{verbatim}
    \end{tcolorbox}
    \caption{Architecture globale à 3 niveaux du projet Edge AI Souss.}
\end{figure}

% =========================================================================
\chapter{Module Embarqué \& Firmware STM32 (Membre 1)}

\section{Spécifications Matérielles}
Le nœud Edge est construit autour de la carte d'évaluation industrielle **STMicroelectronics STM32 Nucleo-F401RE** :
\begin{itemize}
    \item \textbf{Cœur Processeur :} ARM Cortex-M4 32 bits cadencé à 84 MHz.
    \item \textbf{Unité de Calcul Flottant (FPU) :} Simple précision matérielle conforme IEEE 754.
    \item \textbf{Mémoire Flash :} 512 Ko.
    \item \textbf{Mémoire SRAM :} 96 Ko.
    \item \textbf{Périphériques Exploités :}
    \begin{itemize}
        \item Canaux ADC 12 bits : Broche **PA0** connectée au capteur capacitif de sol ($V_{in} \in [0, 3.3\,\text{V}]$ numérisée sur $[0, 4095]$).
        \item Ports GPIO : Broche **PA5** pilotant l'optocoupleur du relais de la pompe.
        \item Compteur de cycles d'horloge matériel : Registre \texttt{DWT->CYCCNT} du bloc de débogage ARM Cortex.
    \end{itemize}
\end{itemize}

\section{Modèle d'Inférence C Pur (\texttt{tinyml\_model\_souss.h})}
Afin de proscrire toute utilisation de mémoire dynamique (`malloc`), le modèle entraîné a été transcrit en code C ANSI hautement optimisé par le compilateur GCC :

\begin{lstlisting}[language=C,caption={Modèle TinyML compilé pour STM32 Cortex-M4.}]
#ifndef TINYML_MODEL_SOUSS_H
#define TINYML_MODEL_SOUSS_H

#include <stdint.h>

/**
 * @brief Inférence TinyML ultra-rapide en C ANSI pur (Arbre de Décision Souss).
 * @param temp_c Température ambiante en °C
 * @param humidity Humidité relative de l'air en %
 * @param et0_mm Évapotranspiration de référence FAO-56 en mm/jour
 * @param soil_moist_raw Valeur brute ADC 12 bits du capteur de sol [0-4095]
 * @return 1 si arrosage nécessaire, 0 sinon
 */
static inline uint8_t tinyml_predict_irrigation(float temp_c, float humidity, 
                                                float et0_mm, float soil_moist_raw) {
    if (soil_moist_raw <= 2200.0f) {
        if (et0_mm > 3.2f || temp_c > 28.0f) {
            return 1; // ARROSAGE ACTIF
        }
        return (humidity < 40.0f) ? 1 : 0;
    } else {
        if (et0_mm > 7.0f && humidity < 25.0f) {
            return 1; // DÉROGATION CHERGUI MAJEUR
        }
        return 0; // REPOS HYDRIQUE
    }
}
#endif
\end{lstlisting}

\section{Métriques Matérielles Officielles (Compilateur ARM-GCC)}
La compilation du module avec `arm-none-eabi-gcc 14.3` (`-O3 -mcpu=cortex-m4 -mthumb -mfpu=fpv4-sp-d16 -mfloat-abi=hard`) certifie les performances d'excellence suivantes :

\begin{table}[H]
    \centering
    \begin{tabular}{lccc}
        \toprule
        \textbf{Caractéristique} & \textbf{Mesure Expérimentale} & \textbf{Capacité Matérielle} & \textbf{Ratio} \\
        \midrule
        Empreinte Flash (Code \texttt{.text}) & \textbf{100 octets} & 512 Ko & 0.019 \% \\
        Empreinte RAM Statique (\texttt{.data} / \texttt{.bss}) & \textbf{0 octet} & 96 Ko & 0.000 \% \\
        Allocation Dynamique (\texttt{malloc}) & \textbf{0 octet (interdite)} & -- & Sécurité absolue \\
        Latence d'Inférence Moyenne & \textbf{0.183 $\mu$s} & $< 1000\,\mu$s & 15.4 cycles CPU \\
        Latence d'Inférence Maximale (Worst-Case) & \textbf{0.202 $\mu$s} & -- & 17.0 cycles CPU \\
        Fréquence d'Horloge Système & \textbf{84 MHz} & 84 MHz & 100 \% \\
        \bottomrule
    \end{tabular}
    \caption{Métriques d'exécution matérielle sur STM32 Nucleo-F401RE.}
\end{table}

\section{Sécurité Industrielle : Coupe-Circuit Watchdog (Anti-Noyade)}
Un mécanisme d'interruption logicielle surveille la durée continue d'activation de la pompe. Tout cycle d'arrosage est rigoureusement borné à un plafond de **30 secondes**. 

Lors du banc d'essai continu de 731 jours, une panne de capteur a été injectée au jour 200 (forçage d'une lecture de sécheresse extrême $48^\circ\text{C}$ et $ET_0 = 12\,\text{mm/j}$). Le Watchdog a automatiquement coupé le relais de la pompe exactement à 30 secondes, annulant tout risque de sur-arrosage (risque d'inondation de **0.0\%**).

% =========================================================================
\chapter{Data Science, Climatologie \& TinyML (Membre 2)}

\section{Dataset Satellitaire Copernicus ERA5}
L'ensemble de données météorologiques a été extrait du service européen de surveillance climatique **Copernicus ERA5** pour la zone du Souss ($30.42^\circ\text{N}, -8.87^\circ\text{W}$) du 1\textsuperscript{er} janvier 2023 au 31 décembre 2024 (731 journées complètes). Les variables climatiques étudiées sont :
\begin{itemize}
    \item Température de l'air à 2 mètres ($T$ en $^\circ\text{C}$).
    \item Humidité relative ($RH$ en \%).
    \item Vitesse du vent normalisée à 2 mètres ($u_2$ en m/s).
    \item Rayonnement solaire global incident ($R_s$ en $\text{MJ/m}^2/\text{jour}$).
    \item Précipitations journalières ($P$ en mm).
\end{itemize}

\section{Modélisation de l'Évapotranspiration de Référence FAO-56}
La demande en eau de l'atmosphère est régie par l'équation internationale standardisée de **Penman-Monteith (FAO-56)** :
\begin{equation}
    ET_0 = \frac{0.408 \Delta (R_n - G) + \gamma \frac{900}{T + 273} u_2 (e_s - e_a)}{\Delta + \gamma (1 + 0.34 u_2)}
\end{equation}
où $R_n$ est le rayonnement net à la surface de la culture, $G$ la densité du flux thermique du sol, $T$ la température moyenne, $u_2$ la vitesse du vent à 2 m, $(e_s - e_a)$ le déficit de pression de vapeur saturante, $\Delta$ la pente de la courbe de pression de vapeur saturante et $\gamma$ la constante psychrométrique.

Le besoin net d'irrigation est ensuite évalué agronomiquement par :
\begin{equation}
    I_{net} = \max\left(0,\, K_c \times ET_0 - P_{eff}\right)
\end{equation}
avec un coefficient cultural maraîcher $K_c = 0.85$ et $P_{eff}$ la pluie efficace.

\section{Benchmarking Comparatif Multi-Modèles}
Trois algorithmes de classification supervisée ont été entraînés et validés sur une partition 75\% apprentissage / 25\% test :

\begin{table}[H]
    \centering
    \begin{tabular}{lcccc}
        \toprule
        \textbf{Modèle Évalué} & \textbf{Accuracy} & \textbf{Sensibilité (Recall)} & \textbf{F1-Score} & \textbf{Poids Mémoire} \\
        \midrule
        \textbf{Arbre de Décision (TinyML)} & \textbf{96.72 \%} & \textbf{98.3 \%} & \textbf{0.98} & \textbf{100 octets} \\
        Forêt Aléatoire (100 arbres) & 97.27 \% & 98.3 \% & 0.98 & $\approx 85$ Ko \\
        Perceptron Multicouche (MLP) & 97.27 \% & 98.3 \% & 0.98 & $\approx 140$ Ko \\
        \bottomrule
    \end{tabular}
    \caption{Comparaison des modèles d'intelligence artificielle.}
\end{table}

\textbf{Justification du Choix de l'Arbre de Décision :}
L'arbre de décision atteint une précision quasi-identique aux modèles ensemblistes et réseaux de neurones (96.72\% vs 97.27\%), tout en réduisant l'empreinte mémoire d'un facteur 1000 et en éliminant toute dépendance à un moteur d'inférence lourd (tel que TensorFlow Lite for Microcontrollers).

% =========================================================================
\chapter{Fog Computing, Supervision \& Benchmark (Membre 3)}

\section{Architecture Fog \& Télémétrie MQTT}
La passerelle de ferme assure la médiation des échanges via le protocole ouvert **MQTT** sur le port 1883 :
\begin{table}[H]
    \centering
    \begin{tabular}{llp{7.5cm}}
        \toprule
        \textbf{Topic MQTT} & \textbf{Sens} & \textbf{Payload / Données Échangées} \\
        \midrule
        \texttt{souss/sensors} & Simulateur $\rightarrow$ Broker & JSON : Horodatage, $T$, $RH$, Vent, $ET_0$, Sol \\
        \texttt{souss/pump} & Bridge $\rightarrow$ Simulateur & État binaire de commande de la pompe (\texttt{ON} / \texttt{OFF}) \\
        \texttt{souss/stm32/log} & Bridge $\rightarrow$ Broker & Télémétrie brute et métriques de latence de la carte \\
        \bottomrule
    \end{tabular}
    \caption{Topologie des flux de messages MQTT.}
\end{table}

\section{Tableau de Bord Node-RED 2.0}
Un flow complet \textit{Node-RED 2.0} a été configuré (\texttt{flow\_souss\_export.json}). Il intègre :
\begin{itemize}
    \item Des jauges visuelles pour la surveillance thermique et l'humidité de l'air.
    \item Des graphes historiques en temps réel pour l'évapotranspiration $ET_0$ et la tension d'eau dans le sol.
    \item Un témoin lumineux d'état de la motopompe couplé à un commutateur d'urgence pour la reprise manuelle.
    \item L'enregistrement chronologique automatique des décisions dans \texttt{decisions.csv}.
\end{itemize}

\section{Benchmark Écologique d'Économie d'Eau}
Deux campagnes de simulation ont été menées pour évaluer le gain hydrique :
\begin{enumerate}
    \item \textbf{Campagne Saisonnière (90 jours, \texttt{benchmark\_eau.py}) :}
    Prise en compte des pertes d'application liées au vent et à l'évaporation diurne. L'arrosage ciblé à l'aube (03h--07h) piloté par l'Edge AI permet d'économiser **15.0\%** du volume d'eau pompé (5 120 $\text{m}^3/\text{ha}$ contre 6 020 $\text{m}^3/\text{ha}$ pour la minuterie classique).
    \item \textbf{Campagne Annuelle Continue (365 jours réels du Souss) :}
    Sous le régime complet du climat de Taroudant, la consommation annuelle mesurée passe de **4 305 Litres** (système à calendrier fixe) à **2 914 Litres** (système Edge AI STM32), réalisant une **économie d'eau de 32.3\%** (1 391 Litres préservés par secteur testé).
\end{enumerate}

% =========================================================================
\chapter{Banc d'Épreuve des 5 Scénarios Climatiques Réels}

Le firmware C pur exécuté sur le processeur ARM Cortex-M4 a été soumis aux cinq conditions météorologiques réelles représentatives de la vallée du Souss :

\begin{table}[H]
    \centering
    \begin{tabular}{cllccc}
        \toprule
        \textbf{N°} & \textbf{Scénario Climatique du Souss} & \textbf{Conditions Météorologiques} & \textbf{Décision} & \textbf{Latence} & \textbf{Statut} \\
        \midrule
        1 & Canicule d'Été (Taroudant) & $T=38.5^\circ\text{C}, RH=20\%, ET_0=7.8$ & POMPE ON & 0.202 $\mu$s & Validé \\
        2 & Journée d'Hiver Humide & $T=16.0^\circ\text{C}, RH=85\%, ET_0=1.8$ & POMPE OFF & 0.155 $\mu$s & Validé \\
        3 & Maraîchage Printanier & $T=25.0^\circ\text{C}, RH=52\%, ET_0=4.6$ & POMPE ON & 0.202 $\mu$s & Validé \\
        4 & Épisode Extrême de Chergui & $T=42.0^\circ\text{C}, RH=12\%, ET_0=9.2$ & POMPE ON & 0.202 $\mu$s & Validé \\
        5 & Automne Tempéré & $T=19.5^\circ\text{C}, RH=68\%, ET_0=2.3$ & POMPE OFF & 0.155 $\mu$s & Validé \\
        \bottomrule
    \end{tabular}
    \caption{Résultats d'exécution du banc de test des scénarios climatiques du Souss.}
\end{table}

\textbf{Taux de conformité : 5 / 5 scénarios validés avec succès (100.0\% de conformité agronomique)}.

% =========================================================================
\chapter{Conclusion \& Perspectives}

Le projet a démontré avec rigueur la viabilité technique, écologique et économique de l'**Edge AI frugale** au service de l'irrigation agricole de précision dans les zones arides du Royaume du Maroc.

\section{Bilan des Résultats Clés}
\begin{itemize}
    \item \textbf{Autonomie Totale :} Décision d'arrosage déterministe en locale sur STM32 en **0.183 microseconde**, immunisée contre les pannes d'Internet.
    \item \textbf{Frugalité Matérielle :} **100 octets de Flash**, aucune utilisation de mémoire dynamique, adaptable sur des microcontrôleurs à très bas coût (< 2 \$).
    \item \textbf{Sécurité Industrielle :} Zéro risque d'inondation (0.0\%) garanti par le Watchdog 30s.
    \item \textbf{Sobriété Hydrique :} Économie d'eau démontrée de **15.0\% à 32.3\%** selon le mode d'exploitation et la saison.
\end{itemize}

\section{Perspectives de Déploiement}
\begin{enumerate}
    \item \textbf{Auto-Alimentation Photovoltaïque :} Intégration d'un module de récupération d'énergie (\textit{Energy Harvesting}) avec un mini-panneau solaire 5W pour une autonomie énergétique illimitée.
    \item \textbf{Réseau Sans-Fil LoRa P2P :} Maillage radio bidirectionnel longue portée entre parcelles sans nécessiter de passerelle 4G sur chaque bloc.
    \item \textbf{Modulation Phénologique :} Ajustement automatique du coefficient $K_c$ au cours du cycle de vie des cultures maraîchères (stade végétatif, floraison, récolte).
\end{enumerate}

\end{document}
